\section{Metodologia}

Conforme apresentado na introdução, este estudo toma como referência o trabalho ``Sistema Embarcado para Reconhecimento de Palavras usando ESP32-S3 e TinyML''~\cite{carlos2025}.      
Todavia, em vez da plataforma Edge Impulse, utilizou-se o PyTorch, mantendo rigorosamente o mesmo conjunto de dados e as mesmas configurações de pré-processamento, arquitetura e      
treinamento adotadas na plataforma original.

\subsection{Conjunto de Dados e Extração de Características (MFCC)}

O conjunto de dados totaliza 1.302~áudios, sendo 1.050 destinados ao desenvolvimento do modelo e 252 mantidos isolados para teste. Cada áudio foi rotulado em uma de seis classes      
(\textit{zero}, \textit{um}, \textit{dois}, \textit{três}, \textit{ruído} e \textit{desconhecido}) e padronizado temporalmente em 16~kHz com duração de 1~segundo (16.000~amostras),   
aplicando-se corte ou preenchimento com zeros (\textit{zero-padding}). Durante a etapa de desenvolvimento, as 1.050~amostras foram divididas em 80\% para treinamento (840~amostras) e 
20\% para validação (210~amostras)~\cite{carlos2025}.

Para extrair as características acústicas utilizadas no reconhecimento de voz, empregaram-se os Coeficientes Mel-Cepstrais (MFCC), técnica que modela a resposta não linear da audição 
humana --- mais sensível às baixas frequências do que às altas ---, seguida de normalização cepstral por média e variância (CMVN)~\cite{carlos2025,pereira2023}. Os parâmetros adotados
reproduzem integralmente a configuração do Edge Impulse e estão resumidos na Tabela~\ref{tab:mfcc}.

\begin{table}[htbp]
\centering
\caption{Parâmetros de extração MFCC e normalização CMVN.}
\label{tab:mfcc}
\begin{tabular}{lc} 
\toprule
\textbf{Parâmetro} & \textbf{Valor} \\ 
\midrule
Coeficientes cepstrais (MFCC) & 13 \\ 
Comprimento e passo do quadro & 20~ms (320 amostras) \\
Número de filtros Mel & 32 \\
Tamanho da FFT & 512~pontos \\
Frequência mínima & 10~Hz \\
Coeficiente de pré-ênfase & 0,98 \\
Janela de normalização CMVN & 101~quadros \\
\bottomrule
\end{tabular}
\end{table}

Ao final do pré-processamento, cada áudio gera uma matriz de entrada de dimensões $(13, 50)$, correspondendo a 13~coeficientes ao longo de 50~quadros temporais. No conjunto de
treinamento, os valores normalizados apresentaram média $-0{,}000$, desvio padrão $0{,}999$ e extremos assimétricos entre $-6{,}362$ e $+5{,}051$.

\subsection{Arquitetura da Rede Neural (1D-CNN)}

A rede neural convolucional unidimensional (1D-CNN) adotada neste trabalho reproduz a arquitetura empregada no Edge Impulse~\cite{carlos2025}, totalizando 1.974~parâmetros treináveis
distribuídos conforme detalhado na Tabela~\ref{tab:arquitetura}.                                                                                                                       

\begin{table}[htbp]
\centering
\setlength{\tabcolsep}{3.5pt}
\caption{Arquitetura da rede 1D-CNN (\textit{KWS\_CNN}) e número de parâmetros.}
\label{tab:arquitetura}
\begin{tabular}{lcc}   
\toprule            
\textbf{Camada} & \textbf{Saída} & \textbf{Parâmetros} \\
\midrule
Entrada MFCC & $(13, 50)$ & --- \\
\texttt{Conv1d} ($13 \to 8$, $k=3$, $p=1$) + ReLU & $(8, 50)$ & 320 \\
\texttt{MaxPool1d} ($k=2$, \textit{ceil}) + \texttt{Dropout} ($0{,}25$) & $(8, 25)$ & --- \\
\texttt{Conv1d} ($8 \to 16$, $k=3$, $p=1$) + ReLU & $(16, 25)$ & 400 \\
\texttt{MaxPool1d} ($k=2$, \textit{ceil}) + \texttt{Dropout} ($0{,}25$) & $(16, 13)$ & --- \\
\texttt{Flatten} + \texttt{Linear} ($208 \to 6$) & $(6)$ & 1.254 \\
\midrule
\textbf{Total} & --- & \textbf{1.974} \\
\bottomrule
\end{tabular}
\end{table}

\begin{figure}[htbp]
\centering
\resizebox{\columnwidth}{!}{\begin{tikzpicture}[
    font=\scriptsize,
    >=Stealth,
    node distance=0.24cm and 0.35cm,
    layer/.style={
        rectangle, rounded corners=2pt, draw=black!75, fill=gray!8,
        minimum width=3.9cm, minimum height=0.44cm, align=center, inner sep=2pt
    },
    quant/.style={
        rectangle, rounded corners=2pt, draw=blue!70!black, fill=blue!10,
        thick, minimum width=1.45cm, minimum height=0.40cm, align=center,
        font=\scriptsize\bfseries, inner sep=2pt
    },
    weight/.style={
        rectangle, rounded corners=2pt, draw=black!60, fill=white,
        minimum width=1.15cm, minimum height=0.40cm, align=center, inner sep=2pt
    },
    io/.style={
        rectangle, rounded corners=2pt, draw=black!80, fill=gray!18,
        minimum width=3.9cm, minimum height=0.44cm, align=center,
        font=\scriptsize\bfseries, inner sep=2pt
    }
]

% Fluxo principal (coluna esquerda)
\node[io]    (inp)   {Entrada MFCC $(13, 50)$};
\node[quant, below=of inp, minimum width=2.4cm] (qin) {Quantizador \texttt{in}};
\node[layer, below=of qin]  (conv1) {\texttt{Conv1d} ($13 \to 8$, $k=3$) + \texttt{ReLU}};
\node[quant, below=of conv1, minimum width=2.4cm] (qa1) {Quantizador \texttt{a1}};
\node[layer, below=of qa1]  (pool1) {\texttt{MaxPool1d} ($k=2$) + \texttt{Dropout} ($0{,}25$)};
\node[layer, below=of pool1] (conv2) {\texttt{Conv1d} ($8 \to 16$, $k=3$) + \texttt{ReLU}};
\node[quant, below=of conv2, minimum width=2.4cm] (qa2) {Quantizador \texttt{a2}};
\node[layer, below=of qa2]  (pool2) {\texttt{MaxPool1d} ($k=2$) + \texttt{Dropout} ($0{,}25$)};
\node[layer, below=of pool2] (flat)  {\texttt{Flatten} ($208$)};
\node[layer, below=of flat]  (fc)    {\texttt{Linear} ($208 \to 6$)};
\node[io,    below=of fc]    (out)   {Saída / \textit{Logits} (6 classes)};

% Ramos laterais de pesos
\node[quant, right=0.45cm of conv1] (qw1)  {Quant. \texttt{w1}};
\node[weight, right=0.30cm of qw1]  (w1)   {Pesos $w_1$};

\node[quant, right=0.45cm of conv2] (qw2)  {Quant. \texttt{w2}};
\node[weight, right=0.30cm of qw2]  (w2)   {Pesos $w_2$};

\node[quant, right=0.45cm of fc]    (qwfc) {Quant. \texttt{wfc}};
\node[weight, right=0.30cm of qwfc] (wfc)  {Pesos $w_{fc}$};

% Conexoes principais
\draw[->] (inp)   -- (qin);
\draw[->] (qin)   -- (conv1);
\draw[->] (conv1) -- (qa1);
\draw[->] (qa1)   -- (pool1);
\draw[->] (pool1) -- (conv2);
\draw[->] (conv2) -- (qa2);
\draw[->] (qa2)   -- (pool2);
\draw[->] (pool2) -- (flat);
\draw[->] (flat)  -- (fc);
\draw[->] (fc)    -- (out);

% Conexoes dos pesos
\draw[->] (w1)   -- (qw1);
\draw[->] (qw1)  -- (conv1);
\draw[->] (w2)   -- (qw2);
\draw[->] (qw2)  -- (conv2);
\draw[->] (wfc)  -- (qwfc);
\draw[->] (qwfc) -- (fc);

\end{tikzpicture}
}
\caption{Arquitetura \textit{KWS\_CNN} com os seis pontos de quantização simulada (\texttt{in}, \texttt{w1}, \texttt{a1}, \texttt{w2}, \texttt{a2} e \texttt{wfc}); os vieses (\textit{biases}) permanecem em FP32.}
\label{fig:arquitetura}
\end{figure}

Para avaliar as estratégias de quantização, a quantização simulada foi inserida em seis pontos da rede: na entrada MFCC (\texttt{in}), nos pesos das camadas convolucionais e da camada
densa (\texttt{w1}, \texttt{w2} e \texttt{wfc}) e nas ativações pós-ReLU (\texttt{a1} e \texttt{a2}), conforme ilustrado na Figura~\ref{fig:arquitetura}. A entrada \texttt{in} utiliza
representação com sinal e simétrica ($\beta = 0$) no PTQ, QAT Padrão e LSQ, passando para representação sem sinal com $\beta$ treinável no LSQ+ e no LSQ+ Híbrido. Os pesos    
(\texttt{w1}, \texttt{w2} e \texttt{wfc}) mantêm quantização com sinal e simétrica ($\beta = 0$) em todas as técnicas. Já as ativações pós-ReLU (\texttt{a1} e \texttt{a2}) adotam 
representação sem sinal em todos os métodos, mantendo $\beta = 0$ no PTQ, QAT Padrão, LSQ e LSQ+ Híbrido, e com $\beta$ treinável apenas no LSQ+ padrão. Os vieses (\textit{biases}) 
das camadas permanecem em ponto flutuante (FP32)~\cite{nagel2021}.

\subsection{Protocolo Experimental e Métricas de Avaliação}
% Treinamento baseline FP32 (300 epocas, Adam, lr=0,005), fine-tuning QAT (100 epocas, Adam, lr=1e-3, CosineAnnealingLR, seed 42), precisoes W8A8, W4A4 e W2A2, e metricas (Acuracia, Loss, ROC AUC e F1-Score ponderado).
